\chapter{Mar.~5 --- Coherent Sheaves on Projective Varieties}

\section{Coherent Sheaves on Projective Varieties}

\begin{remark}
  Let $X \subseteq \PP^n$ be a projective
  variety (i.e. a closed subvariety), and
  let \[S = S(X) = k[x_0, \dots, x_n] / I(X),\]
  which is a graded ring.
  Recall the notion of a graded $S$-module
  (Definition \ref{def:graded-module}).
\end{remark}

\begin{example}
  Trivially, $S$ is a graded $S$-module.
\end{example}

\begin{example}
  If $M$ is a graded $S$-module, then
  $M(m)$ is the graded $S$-module with
  \[
    M(m)_d = M_{d + m}.
  \]
\end{example}

\begin{remark}
  $\mathrm{Gr}\text{-}\mathrm{Mod}_S$
  is an abelian category.
\end{remark}

\begin{remark}
  We want a correspondence
  \[
    \begin{tikzcd}
    \text{(finitely generated) graded $S$-modules}
    \ar[r, shift left = 1] & \text{(coherent) quasicoherent sheaves on $X$} \ar[l, shift left = 1]
    \end{tikzcd}
  \]
  To begin with, we want $S$ to correspond
  to $\OO_X$.
\end{remark}

\begin{example}
  Let $X = \PP^n$, then $S = k[x_0, \dots, x_n]$.
  Note that
  \[
    \OO_{\PP^n}(D(x_i))
    = k[x_0 / x_i, \dots, x_n / x_i]
    = ((k[x_0, \dots, x_n])_{x_i})_0
    = \left\{
      \frac{f}{x_i^d} : f \in k[x_0, \dots, x_n]_d
    \right\}
  \]
  Note that $k[x_0, \dots, x_n]_{x_i}$
  is naturally graded. More generally,
  for $f \in S^+$ of homogeneous
  degree $d$, let
  \[
    S_{(f)}
    := (S_f)_0
    = \text{degree $0$ part of $S_f$}
    = \left\{
      \frac{g}{f^i} \in S_f : g \in S_{di},\,
      i \in \N
    \right\}.
  \]
  Then there is a homomorphism
  \begin{align*}
    \phi : S_{(f)} &\longrightarrow \OO_X(D(f)) \\
    g / f^i &\longmapsto (x \mapsto g(x) / f^i(x)),
  \end{align*}
  which we claim is an isomorphism.
  To see this, consider $Y = C(X) \subseteq \Affine^{n + 1}$, which
  satisfies $A(Y) = S$.
  Now observe that we have a commutative
  diagram
  \[
    \begin{tikzcd}
      S_{(f)} \ar[r, "\phi"] \ar[d, hook] & \OO_X(D_X(f)) \ar[d, hook] \\
      S_f \ar[r, "\cong"] & \OO_Y(D_Y(f))
    \end{tikzcd}
  \]
  Thus $\phi$ is injective. Furthermore,
  if $\psi \in \OO_X(D(f))$, we get
  \[
    \psi = g / f^i, \quad g \in S(X).
  \]
  Then $g(\lambda x) = \lambda^{d i} g(x)$
  for all $x \in X$ such that $g(x) \ne 0$.
  So $g \in S_{di}$ (assuming the previous
  condition holds on a dense subset of
  $X$, e.g. if $g$ is nonzero on every
  irreducible component).
\end{example}

\begin{remark}
  Fix $M \in \mathrm{Gr}\text{-}\mathrm{Mod}_S$
  and $f \in S^+$ homogeneous of degree $d$.
  We want an $\OO_X$-module. Note that
  \[
    M_{(f)}
    := \text{degree $0$ part of $M_f$}
    = \left\{
      \frac{m}{f^i} \in M_f : m \in M_{di},\,
      i \in \N
    \right\}
  \]
  is an $S_{(f)}$-module. Then we can set
  \begin{itemize}
    \item $\widetilde{M}(D(f)) = M_{(f)}$;
    \item For $f, g \in S^+$ homogeneous,
      we have natural maps
      $M_{(f)} \to M_{(fg)}$ and $M_{(g)} \to M_{(fg)}$,
      which we can view as the restriction
      maps
      $\widetilde{M}(D(f)) \to \widetilde{M}(D(fg))$ and
      $\widetilde{M}(D(g)) \to \widetilde{M}(D(fg))$.
  \end{itemize}
\end{remark}
Then we claim
that $\widetilde{M}$ is a
sheaf of $\OO_X$-modules on the
basis $\{D(f) : f \in S^+ \text{ homogeneous}\}$,
hence it induces $\widetilde{M} \in \mathrm{Mod}_{\OO_X}$. The
proof is similar to the affine case.

\begin{prop}
  $\widetilde{S(m)}$ is an invertible
  sheaf. In particular, it is quasicoherent.
\end{prop}

\begin{proof}
  Write $X = D(x_0) \cup \cdots \cup D(x_n)$.
  Now we have an isomorphism
  \begin{align*}
    \widetilde{S(m)}(D(x_i))
    = (S_{x_i})_m
    &\overset{\cong}{\longrightarrow}
    (S_{x_i})_0 \\
    r &\longmapsto x_i^{-m} r \\
    x_i^m r &\mathrel{\reflectbox{\ensuremath{\longmapsto}}} r.
  \end{align*}
  Note that $D(x_i)$ has a basis by
  $\{D(x_i g) : g \in S^+ \text{ homogeneous}\}$,
  and
  \begin{align*}
    \widetilde{S(m)}(D(x_i g))
    = (S_{x_i g})_m
    &\overset{\cong}{\longrightarrow}
    (S_{x_i g})_0 \\
    r &\longmapsto x_i^{-m} r
  \end{align*}
  As this isomorphism is compatible
  with the previous one, we get that
  $\widetilde{S(m)}|_{D(x_i)} \cong \OO_{D(x_i)}$.
  So we see that
  $\widetilde{S(m)}$ is an invertible sheaf.
\end{proof}

\begin{remark}
  We have $\OO_X(m) \cong \widetilde{S(m)}$.
  To see this, note that we have an
  isomorphism
  \begin{align*}
    \alpha_i : \widetilde{S(m)}|_{D(x_i)}
    &\longrightarrow
    \OO_{D(x_i)} \\
    x_i^m  &\longmapsto 1.
  \end{align*}
  Now observe that we have the composition
  \[
    \begin{tikzcd}[row sep=tiny]
      \OO_{D(x_i x_j)} \ar[r, "\alpha_j^{-1}"]
      & \widetilde{S(m)}|_{D(x_i x_j)} \ar[r, "\alpha_i"]
      & \OO_{D(x_i x_j)} \\
      1 \ar[r, mapsto] & x_j^m = (x_i^m)(x_j / x_i)^m
      \ar[r, mapsto] & (x_j / x_i)^m
    \end{tikzcd}
  \]
  so the transition functions are
  $g_{i, j} = (x_j / x_i)^m$. Thus
  we see that $\widetilde{S(m)} \cong \OO_X(m)$.
\end{remark}

\begin{remark}
  Consider the map
  $\mathrm{Gr}\text{-}\mathrm{Mod}_S \to \mathrm{Mod}_{\OO_X}$
  given by the construction
  $M \mapsto \widetilde{M}$ above. Then:
  \begin{enumerate}
    \item This is a functor.
    \item This is exact (localization
      is exact and the $D(f)$ form a basis).
    \item If $M_i \in \mathrm{Gr}\text{-}\mathrm{Mod}_S$
      for $i \in I$, then
      $\widetilde{\bigoplus_{i \in I} M_i} \cong \bigoplus_{i \in I} \widetilde{M}_i$
      (localization commutes with
      direct sums).
  \end{enumerate}
\end{remark}

\begin{prop}
  For $M \in \mathrm{Gr}\text{-}\mathrm{Mod}_S$,
  $\widetilde{M}$ is quasicoherent.
  Furthermore, if $M$ is finitely generated,
  then $\widetilde{M}$ is coherent.
\end{prop}

\begin{proof}
  We have an exact sequence
  \[
    \begin{tikzcd}
      \bigoplus_i S(-a_i) \ar[r]
      & \bigoplus_j S(-b_j) \ar[r]
      & M \ar[r] & 0
    \end{tikzcd} \tag{$*$}
  \]
  Applying $M \mapsto \widetilde{M}$
  (which is exact and commutes with direct sums),
  we get an exact sequence
  \[
    \begin{tikzcd}
      \bigoplus_i \OO_X(-a_i) \ar[r]
      & \bigoplus_j \OO_X(-b_j) \ar[r]
      & \widetilde{M} \ar[r] & 0
    \end{tikzcd}
  \]
  This gives that $\widetilde{M}$ is
  quasicoherent. If $M$ is finitely
  generated, then we can choose the
  direct sums in $(*)$ to be finite, hence
  $\widetilde{M}$ is coherent in this case.
\end{proof}

\begin{remark}
  We have $\widetilde{M} \otimes \OO_X(m) \cong \widetilde{M(m)}$.
  To see this, we can note that
  \[
    (\widetilde{M} \otimes \OO_X(m))(D(x_i))
    = (M_{x_i})_0 \otimes_{(S_{x_i})_0} (S_{x_i})_m
    \cong (M_{x_i})_m
    = \widetilde{M(m)}(D(x_i)).
  \]
  Alternatively, shifting the exact sequence
  $\bigoplus_i S(-a_i) \to \bigoplus_j S(-b_j) \to M \to 0$ by $m$
  (which preserves exactness),
  \[
    \begin{tikzcd}
      \bigoplus_i S(m - a_i) \ar[r]
      & \bigoplus_j S(m - b_j) \ar[r]
      & M(m) \ar[r] & 0.
    \end{tikzcd}
  \]
  Applying $M \mapsto \widetilde{M}$, we get
  that
  \[
    \widetilde{M(m)}
    = \coker\Big(\bigoplus_{i} \OO_X(m - a_i) \to \bigoplus_j \OO_X(m - b_j) \Big)
    = \widetilde{M} \otimes \OO_X(m).
  \]
\end{remark}

\begin{remark}
  When is $\widetilde{M} = 0$? Clearly
  this is true if $M = 0$,
  but there are more examples.
  Observe that
  we can write
  $s / x_i^d = x_i^c s / x_i^{d + c} \in M_{x_i}$,
  so if $x_i^c s = 0$ for some $c > 0$, then
  $s / x_i^d$ is zero in $\widetilde{M}(D(x_i))$.

  In particular, if there exists $c > 0$
  such that $M_i = 0$ for all $i \ge c$,
  then $\widetilde{M}(D(x_i)) = 0$, so
  $\widetilde{M} = 0$.
\end{remark}

\begin{corollary}
  If $\alpha : M \to M'$ is a morphism of graded
  $S$-modules and $M_i \to M_i'$ is an
  isomorphism for $i \gg 0$, then
  $\widetilde{M} \to \widetilde{M}'$
  is an isomorphism.
\end{corollary}

\begin{proof}
  Consider the exact sequence of graded
  modules
  \[
    \begin{tikzcd}
      0 \ar[r] & \ker \alpha \ar[r] & M \ar[r, "\alpha"] & M' \ar[r] & \coker \alpha \ar[r] & 0
    \end{tikzcd}
  \]
  Now $\ker \alpha_i = 0 = \coker \alpha_i$
  for $i \gg 0$. So we get an exact sequence
  \[
    \begin{tikzcd}
      0 \ar[r] & 0 \ar[r] & \widetilde{M} \ar[r] & \widetilde{M}' \ar[r] & 0 \ar[r] & 0
    \end{tikzcd}
  \]
  Hence $\widetilde{M} \to \widetilde{M}'$
  is an isomorphism.
\end{proof}

\section{Graded Modules from Coherent Sheaves}

\begin{remark}
  How do we get a functor
  $\Phi : \QCoh_X \to \mathrm{Gr}\text{-}\mathrm{Mod}_S$?
\end{remark}

\begin{example}
   For $X = \PP^n$, we want
   \[
     \Phi(\OO_{\PP^n})
     = k[x_0, \dots, x_n]
     = \bigoplus_{d \in \N} k[x_0, \dots, x_n]_d
     = \bigoplus_{d \in \N} \Gamma(\PP^n, \OO_{\PP^n}(d)).
   \]
\end{example}

\begin{definition}
  For $\mathcal{F} \in \QCoh_X$, set
  (below $\mathcal{F}(m) = \mathcal{F} \otimes \OO_X(m)$)
  \[
    \Gamma_*(X, \mathcal{F}) = \bigoplus_{m \in \Z} \Gamma(X, \mathcal{F}(m)).
  \]
\end{definition}

\begin{remark}
  For $M \in \mathrm{Gr}\text{-}\mathrm{Mod}_S$,
  we have a morphism of abelian groups
  \[
    M \longrightarrow
    \Gamma_*(X, \widetilde{M})
  \]
  given by the natural maps
  $M_i \to \Gamma(X, \widetilde{M(i)}) \subseteq \Gamma_*(X, \widetilde{M})$.
  In particular, we get
  \[
    S \longrightarrow \Gamma_*(X, \OO_X).
  \]
  Using this, we can get a map
  \[
    \begin{tikzcd}
    S_m \times \Gamma(X, \mathcal{F}(n))
    \ar[r] & \Gamma(X, \OO_X(m)) \times \Gamma(X, \mathcal{F}(n)) \ar[d] \\
    & \Gamma(X, \mathcal{F}(m + n))
    \end{tikzcd}
  \]
  which we can think of as a
  multiplication. Thus
  $\Gamma_*(X, \mathcal{F})$ has the
  structure of a graded $S$-module.
\end{remark}

\begin{prop}
  For $\mathcal{F} \in \QCoh_X$,
  there exists a natural isomorphism
  \[
    \mathcal{F}
    \overset{\cong}{\longrightarrow}
    \widetilde{\Gamma_*(X, \mathcal{F})}.
  \]
\end{prop}

\begin{proof}
  See Mustata 11.1.18.
\end{proof}

\begin{remark}
  The functors
  $(\sim, \Gamma_*)$
  form an adjoint pair.
\end{remark}

\begin{corollary}
  If $\mathcal{F} \in \Coh_X$, then
  there exists a finitely generated
  $S$-module $M$ such that
  $\mathcal{F} = \widetilde{M}$.
\end{corollary}

\begin{proof}
  Set $N = \Gamma_*(X, \mathcal{F})$
  and choose $M \subseteq N$ finitely generated
  with $\widetilde{M} \cong \mathcal{F}$.
\end{proof}
